\section{LWE on real numbers}

\anote{Still some problems that we don't have time to fix in the proof, so maybe it is better to remove this part unless someone have some time to spend on it (proof of claim~\ref{cl:real_left_size} and in claim~\ref{cl:real_right_side} the distribution of $\eps$ need to be studied conditioned on $\vec{s}$ ?}

\subsection{Definitions}

We now consider a version of the LWE problem on real numbers. This version does not
involve a modulus $q$ anymore.

Let $\phi$ be some distribution on~$\mathbb{T}$ and~$\vec{s} \in \Z^n$ be an
arbitrary vector.  We define real-$A_{\vec{s},\phi}$ as the
distribution on~$\mathbb{T}^{n+1}$ obtained by sampling a
vector $\vec{a}$ uniformly in~$\mathbb{T}^{n}$, and $e \in \mathbb{T}$
from distribution~$\phi$, and returning~$(\vec{a}, \inner{\vec{a},\vec{s}}+ e \bmod 1)$.

\begin{definition}
  Let~$n \geq 1$ be an integer, $\beta \in (0,1)$ and~$\Upsilon$ be a family of distributions on~$\mathbb{T}$ (resp. a distribution on such a family). The parameters~$n, \beta$ and~$\Upsilon$ are functions of a parameter~$\lambda$ against which complexity is measured. 
  
  The decision \emph{Learning With Error problem on real}, \emph{real-$\text{LWE}_{\Upsilon}(D_{\Z^n, \beta})$}, is as follows: Let $\vec{s} \in \Z^n$ be sampled from $D_{\Z^n, \beta}$ and~$\phi$ be sampled from~$\Upsilon$; The goal is to distinguish between arbitrarily many independent samples from $A_{\vec{s},\phi}$ and the same number of independent samples from $U( \mathbb{T}^{n+1})$.
\end{definition}

By abuse of notation, we also call this problem real-$\text{LWE}_{\Upsilon}(D_{\Z^n, \beta})$ if~$\Upsilon$
is only a set of distributions and $\phi$ is \emph{arbitrarily} chosen in~$\Upsilon$.


\subsection{Reduction from LWE$_{n,q,\nu_{\sqrt{2}\alpha}}(D_{\Z^n, \alpha q})$ to real-$\text{LWE}_{n,\Upsilon_{\beta_2}}(D_{\Z^n, \beta_1})$}


\begin{theorem}\label{th:real}
  Let~$n \geq 1$, $q \in [2, 2^{n^{O(1)}}]$ and $\alpha, \beta_1, \beta_2
  \in (0,1)$, all parameters of the complexity variable~$\lambda$.
%  Let~$\mathcal{D}$ be an arbitrary distribution of $\Z_q^n$ and define~$s_{\max}(\mathcal{D})$ 
%  as any real such that~$\Pr_{\vec{s} \hookleftarrow \mathcal{D}}[ \Vert \vec{s} \Vert \geq s_{\max}(\mathcal{D}) ] \leq \lambda^{- \omega(1)}$.
  
  Assume that $\beta_1 \geq 2 \alpha q $, 
  $\beta_2 \geq \alpha \mult \omega(\sqrt{n \log \lambda} + \log \lambda)$ and~$\alpha q \geq \omega(\sqrt{\log \lambda})$.  
  Then there exists a probabilistic polynomial-time
  reduction from \emph{LWE$_{n,q,\nu_{\sqrt{2}\alpha}}(D_{\Z^n, \alpha q})$}
  to \emph{real-$\text{LWE}_{n,\nu_{\leq \beta_2}}(D_{\Z^n, \beta_1})$}.
\end{theorem}

\begin{proof}
Assume that we have a polynomial-time oracle that solves real-$\text{LWE}_{n,\nu_{\leq \beta_2}}(D_{\Z^n, \beta_1})$ with probability $\lambda^{-O(1)}$. Our inputs are $m$ samples from
either~$U(\Z_q^n \times \mathbb{T})$ or from~$A_{\vec{s},\nu_{ \alpha}}$, 
where the secret~$\vec{s}$ is distributed from $\mathcal{D}= D_{\Z^n,\alpha q}$. With overwhelming probability, we have
$ \Vert \vec{s} \Vert \leq \alpha q \mult \omega(\sqrt{n} + \sqrt{\log \lambda})$, and we assume in the following that the secret is indeed short.
 The reduction~$\mathcal{R}$ is as follows: 
 
 \begin{itemize}
\item[$\bullet$] Sample $\vec{s}_1$ from $D_{\Z^n, \sqrt{\beta_1^2 - \alpha^2 q^2} }$.
\item[$\bullet$] For each sample $(\vec{a},b)$, create the sample $(\vec{a}',b')$ as follows:
\begin{itemize}
\item[$\bullet$] Let $\sigma = \frac{1}{q} \omega (\sqrt{ \log \lambda })$,
define $\vec{a}' = \frac{1}{q} \vec{a} + \nu_{\sigma} (\bmod~1)$ and $b' = b + \inner{\vec{a}', \vec{s}_1}$.
\item[$\bullet$] Return the new sample $(\vec{a}',b')$.
\end{itemize}
\item[$\bullet$] Call the real-$\text{LWE}_{n,\nu_{\leq \beta_2}}(D_{\Z^n, \beta_1})$ oracle on these
  samples, and return its answer.
\end{itemize} 
Assume that~$b \in \mathbb{T}$ is of the form~$\frac{1}{q}
\inner{\vec{a},\vec{s}}+ e$ with~$e \hookleftarrow \nu_{\alpha}$. Then
we can write:
\begin{equation}
\label{eq:real_rewrite_b}
\left(\vec{a'}, b\right) = 
\left(
\vec{a'}, \inner{\vec{a}',\vec{s} + \vec{s}_1} +  \inner{\vec{\bfeps},\vec{s}} + e  \right).  
\end{equation}
We consider the joint distribution of the pair~$(\vec{a}',\vec{\bfeps})$.  

\begin{claim}
\label{cl:real_left_size}
Assume that~$\sigma \geq \frac{1}{q}\omega(\sqrt{\log \lambda})$. Then the residual
  distribution of~$\vec{a}'$ is within negligible statistical distance
  to~$U(\mathbb{T}^n)$, and for any~$\overline{\vec{a}'} \in \mathbb{T}^n$, the
  distribution of~$\vec{\bfeps}$ conditioned on~$\vec{a}' =
  \overline{\vec{a}'}$ is within negligible statistical distance
  to~$D_{ \overline{\vec{a}'} + \frac{1}{q} \Z^n,\sigma}$.
\end{claim}


\begin{proof}
TODO
\end{proof}

We now analyze the distributions of~$\vec{s} + \vec{s}_1$ and, of the term~$\inner{\vec{\bfeps},\vec{s}} + e $ from Eq.~\eqref{eq:real_rewrite_b}, given~$\vec{a}'$.

\begin{claim}
\label{cl:real_right_side}
  If~$\beta_1 \geq 2 \alpha q$ and~$\alpha q \geq \omega ( \sqrt{\log \lambda})$, then the distribution of~$\vec{s} + \vec{s}_1$
is within negligible statistical distance to~$D_{\mathbb{Z}^n,\beta_1} $.

If~$\vec{\bfeps} \hookleftarrow \nu_{\sigma}$ and~$e \hookleftarrow \nu_{\alpha}$, then the
  distribution of~$ \inner{\vec{\bfeps},\vec{s}} + e$ is distributed from~$\nu_{\beta}$ with~$\beta=\sqrt{\sigma^2\|\vec{s}\|^2+ \alpha^2}$.

\end{claim}

\begin{proof}
 If~$\beta_1 \geq 2 \alpha q$ and~$\alpha q \geq \omega ( \sqrt{\log \lambda})$, then~$\sqrt{\beta_1^2 - \alpha^2 q^2} \geq \alpha q \geq \omega ( \sqrt{\log \lambda})$. The conditions of Lemma~\ref{le:discrete_plus_cont_Gauss} hold, as a consequence the distribution of~$\vec{s} + \vec{s}_1$ is within negligible statistical distance to~$D_{\mathbb{Z}^n,\beta_1} $.


In the second part of the Claim, we simply add continuous gaussian distributions.
If~$\vec{\bfeps} \hookleftarrow \nu_{\sigma}$ , then $\inner{\vec{\bfeps},\vec{s}}$ is distributed from 
$\nu_{\sigma\|\vec{s}\|}$. As a consequence, we have that $ \inner{\vec{\bfeps},\vec{s}} + e$ is distributed from~$\nu_{\beta}$ with~$\beta=\sqrt{\sigma^2\|\vec{s}\|^2+ \alpha^2}$.
\end{proof}

Let~$\vec{s}' = \vec{s} + \vec{s}_1$.
At this stage, we have proven that~$\vec{a}'$ is within negligible statistical distance to~$U(\mathbb{T}^n)$, and that~$b' = \inner{\vec{a}',\vec{s}'} + e'$, where the distribution of~$e'$ given~$\vec{a}'$ is within negligible statistical distance
to~$\nu_{\sqrt{\sigma^2\|\vec{s}\|^2+ \alpha^2}}$. We have~$\sigma = \frac{1}{q} \omega (\sqrt{\log \lambda})$ and with overwhelming probability~$ \Vert \vec{s} \Vert \leq \alpha q \mult \omega(\sqrt{n} + \sqrt{\log \lambda})$.
Thus, with the same probability, $\sqrt{\sigma^2\|\vec{s}\|^2+ \alpha^2} \leq \alpha \mult \omega(\sqrt{n \log \lambda} + \log \lambda)$.

The random variable~$\vec{s}'$ has distribution~$D_{\mathbb{Z}^n,\beta_1}$, and the standard deviation of the new error is below~$\beta_2$. This conclude the proof of the Theorem~\ref{th:real}.


\end{proof}

